\section{Background}
\label{sec:background}

\subsection{LLM Agent Reinforcement Learning}

In LLM agent RL, the policy \(\pi_\theta\) is a language model. At each environment turn, the model receives a textual observation and produces a response containing both reasoning text and an action. The environment parses the action, transitions to the next state, and returns a reward. This differs from single-turn RLHF~\citep{ouyang2022training}, preference optimization such as DPO~\citep{rafailov2023dpo}, and reasoning-oriented RL such as DeepSeek-R1~\citep{guo2025deepseekr1}, because the generated tokens are embedded in an external multi-turn trajectory.

The objective is to maximize expected return,
\begin{equation}
J(\theta)=\mathbb{E}_{\tau\sim\pi_\theta}
\left[\sum_{t=0}^{T}\gamma^t r_t\right],
\end{equation}
where \(r_t\) is the turn-level environment reward. This project follows RAGEN in using sparse FrozenLake rewards, a small format penalty, KL regularization against a frozen reference policy, and bi-level credit assignment.

\subsection{PPO and GRPO}

PPO~\citep{schulman2017ppo} uses an old policy \(\pi_{\theta_{\mathrm{old}}}\), an importance ratio \(r_t(\theta)=\pi_\theta(a_t\mid s_t)/\pi_{\theta_{\mathrm{old}}}(a_t\mid s_t)\), and a clipped surrogate objective:
\begin{align}
L^{\mathrm{CLIP}}(\theta)
&= \mathbb{E}_t\left[\min(u_t, v_t)\right], \\
u_t &= r_t(\theta)\hat{A}_t, \\
v_t &= \tilde{r}_t(\theta)\hat{A}_t, \\
\tilde{r}_t(\theta)
&= \mathrm{clip}\bigl(r_t(\theta),1-\epsilon,1+\epsilon\bigr).
\end{align}
The implementation adds a value loss, entropy regularization, and KL regularization. It uses the RAGEN-aligned coefficients \(c_{\mathrm{vf}}=1.0\), \(c_{\mathrm{ent}}=0.001\), \(\beta=0.001\), and \(\epsilon=0.2\). PPO computes advantages with bi-level GAE: a turn-level pass estimates value at environment boundaries, and a token-level pass distributes the turn signal across response tokens.

GRPO~\citep{shao2024deepseekmath} removes the critic. Given a group of trajectories sampled from the same prompt, it converts returns into group-relative z-scores:
\begin{equation}
\hat{A}_i = \frac{R_i-\mu_g}{\sigma_g+\varepsilon},
\qquad
\mu_g=\frac{1}{G}\sum_{j=1}^{G}R_j.
\end{equation}
The resulting scalar advantage is applied to the response tokens of that trajectory. Since positive affine transformations of all rewards in a group leave the z-scores approximately unchanged, GRPO mainly uses within-group ranking and relative separation rather than absolute reward scale.

\subsection{StarPO-S Filtering}

StarPO organizes multi-turn LLM agent RL around textual state, thinking text, actions, and rewards. Each turn requires a think/answer output format, and malformed outputs receive a per-turn penalty of \(-0.1\).

StarPO-S adds variance-based rollout filtering. At each update, \(P\) prompts are sampled and \(K\) trajectories are generated per prompt. For each prompt \(p\), the trainer computes the variance of its \(K\) trajectory returns and keeps only the top \(r\cdot P\) prompts, where \(r\) is \code{variance\_filter\_ratio}. RAGEN reports that moderate filtering reduces echo-trap behavior in small-model PPO-style training. The central question here is whether this mechanism still helps after switching from PPO to GRPO.
